\makeabstract
{
Quantifying biofilm exopolysaccharide expression in Pseudomonas aeruginosa clinical strain CF23
}
{
Mahnoor Mustafa(1), Thomas Murray (2), Pamela Huang (2)
}
{
\textsuperscript{1}Amherst College, \textsuperscript{2}Yale University School of Medicine
}
{
To investigate the biofilm forming capacity of the P. aeruginosa CF23 strain by quantifying the expression of biofilm associated genes and then correlating gene expression data with phenotypic traits such as motility and EPS production assays in order to understand the underlying molecular mechanisms that contribute to the strain’s chronic infection phenotype.
}
{
For phenotypic analysis, colony morphology and EPS production were measured using Congo red plates. Furthermore, motility assays were conducted. Twitching motility was assessed by staining with Coomassie Blue to visualize twitching zones at the agar–petri dish interface. Swimming motility was measured by diameter of migration on LB plates containing 0.3\% agar. 
For gene expression analysis, RNA was reverse-transcribed into cDNA to perform quantitative PCR with primers targeting pelA, pslD, algD, and the housekeeping gene rplU. Amplification was carried out, and product specificity was confirmed by melt curve analysis. Amplification products were further validated by gel electrophoresis.
}
{
CF23 (clinical strain) and PA15 (Pel constitutive mutant) displayed stronger Congo red binding when compared to Pel deficient strains (PA25 and PA26). Pel-deficient strains (PA25, PA26) exhibited larger twitching and swimming zones compared to high-EPS strains (CF23, PA15), suggesting that reduced EPS production enhances Type IV pili- and flagella-mediated motility. CF23 expressed pelA, pslD, and algD, confirming robust production of all three EPS types at the transcriptional level.

}
{
These findings suggest that Pseudomonas aeruginosa CF23 strain possesses strong biofilm forming capacity, driven elevated EPS production of Pel, Psl and alginate. The correlation between gene expression data and phenotypic traits such as decreased motility but strong biofilm may suggest how the strain has adapted to chronic infection conditions.
}

\newpage
